\documentclass{beamer}
\usepackage[utf8]{inputenc}
\usepackage[spanish]{babel}
\usepackage{listings}
\usepackage{xcolor}
\usepackage{tikz}
\usepackage{tcolorbox}
\usepackage{fontawesome5}

% Configuración de colores - Gruvbox Light con marrón principal
\definecolor{bg0}{HTML}{fbf1c7}
\definecolor{bg1}{HTML}{ebdbb2}
\definecolor{bg2}{HTML}{d5c4a1}
\definecolor{bg3}{HTML}{bdae93}
\definecolor{fg0}{HTML}{282828}
\definecolor{fg1}{HTML}{3c3836}
\definecolor{fg2}{HTML}{504945}
\definecolor{aqua}{HTML}{689d6a}
\definecolor{aquabright}{HTML}{427b58}
\definecolor{blue}{HTML}{458588}
\definecolor{red}{HTML}{cc241d}
\definecolor{orange}{HTML}{d65d0e}
\definecolor{yellow}{HTML}{d79921}
\definecolor{green}{HTML}{98971a}
\definecolor{purple}{HTML}{b16286}
\definecolor{brown}{HTML}{AB886D}
\definecolor{brownbright}{HTML}{8B6849}

% Colores para código
\definecolor{codegreen}{HTML}{98971a}
\definecolor{codegray}{HTML}{7c6f64}
\definecolor{codepurple}{HTML}{b16286}
\definecolor{backcolour}{HTML}{f9f5d7}

% Configuración de código
\lstdefinestyle{pythonstyle}{
    backgroundcolor=\color{backcolour},   
    commentstyle=\color{codegreen},
    keywordstyle=\color{brownbright}\bfseries,
    numberstyle=\tiny\color{codegray},
    stringstyle=\color{codepurple},
    basicstyle=\ttfamily\scriptsize\color{fg0},
    breakatwhitespace=false,         
    breaklines=true,                 
    captionpos=b,                    
    keepspaces=true,                 
    numbers=left,                    
    numbersep=2pt,                  
    showspaces=false,                
    showstringspaces=false,
    showtabs=false,                  
    tabsize=2,
    language=Python,
    frame=single,
    rulecolor=\color{brown},
    xleftmargin=6pt,
    xrightmargin=6pt,
    inputencoding=utf8,
    extendedchars=true,
    literate=
        {á}{{\'a}}1 {é}{{\'e}}1 {í}{{\'i}}1 {ó}{{\'o}}1 {ú}{{\'u}}1
        {Á}{{\'A}}1 {É}{{\'E}}1 {Í}{{\'I}}1 {Ó}{{\'O}}1 {Ú}{{\'U}}1
        {ñ}{{\~n}}1 {Ñ}{{\~N}}1
        {¿}{{?`}}1 {¡}{{!`}}1
}

\lstset{style=pythonstyle}

% Tema personalizado con marrón principal
\usetheme{Madrid}
\setbeamercolor{structure}{fg=brown}
\setbeamercolor{palette primary}{bg=brown,fg=bg0}
\setbeamercolor{palette secondary}{bg=brownbright,fg=bg0}
\setbeamercolor{palette tertiary}{bg=fg1,fg=bg0}
\setbeamercolor{palette quaternary}{bg=brown,fg=bg0}
\setbeamercolor{block title}{bg=brown,fg=bg0}
\setbeamercolor{block body}{bg=bg1,fg=fg0}
\setbeamercolor{frametitle}{bg=brown,fg=bg0}
\setbeamercolor{title}{fg=fg0}
\setbeamercolor{subtitle}{fg=fg1}
\setbeamercolor{normal text}{fg=fg0,bg=bg0}
\setbeamercolor{item}{fg=brown}
\setbeamercolor{section in toc}{fg=brown}

% Cajas personalizadas con marrón
\newtcolorbox{conceptbox}[1]{
    colback=bg1,
    colframe=brown,
    fonttitle=\bfseries\small,
    title=#1,
    coltitle=fg0,
    left=3pt,
    right=3pt,
    top=3pt,
    bottom=3pt,
    boxsep=2pt
}

\newtcolorbox{notebox}{
    colback=yellow!15,
    colframe=orange,
    leftrule=2mm,
    rightrule=0mm,
    toprule=0mm,
    bottomrule=0mm,
    arc=0mm,
    left=3pt,
    right=3pt,
    top=2pt,
    bottom=2pt
}

\title{\textbf{Introducción a la Programación con Python}}
\subtitle{Semana 7: Expresiones Regulares}
\author{$\nabla$Nablabs}
\date{}
\begin{document}

\begin{frame}
\titlepage
\end{frame}

\begin{frame}
\frametitle{Contenido}
\tableofcontents
\end{frame}

\section{¿Qué son las expresiones regulares?}

\begin{frame}{Expresiones Regulares (Regex) \faSearch}
\begin{conceptbox}{¿Qué son?}
\small
\begin{itemize}
    \item[\faCode] Patrones de texto para buscar, validar y manipular strings
    \item[\faMagic] Lenguaje especializado para procesamiento de texto
    \item[\faCheckCircle] Validar formatos (emails, teléfonos, URLs, etc.)
    \item[\faExchangeAlt] Extraer información específica de textos
    \item[\faEdit] Reemplazar patrones complejos
\end{itemize}
\end{conceptbox}

\vspace{0.2cm}
\begin{center}
\large\textcolor{brown}{\faLightbulb\ ¡Superpoder para trabajar con texto!}
\end{center}
\end{frame}

\begin{frame}{¿Por qué usar regex? \faQuestion}
\begin{columns}[T]
\column{0.48\textwidth}
\textbf{Sin regex:} \faTimesCircle
\small
\begin{itemize}
    \item Código largo y complejo
    \item Muchas condiciones
    \item Difícil de mantener
    \item Validaciones limitadas
    \item Procesamiento lento
\end{itemize}

\column{0.48\textwidth}
\textbf{Con regex:} \faCheckCircle
\small
\begin{itemize}
    \item Código conciso
    \item Patrones expresivos
    \item Fácil de actualizar
    \item Validaciones potentes
    \item Procesamiento eficiente
\end{itemize}
\end{columns}

\vspace{0.3cm}
\begin{center}
\textcolor{brown}{\faInfoCircle\ Regex es universal: Python, JavaScript, Java, etc.}
\end{center}
\end{frame}

\section{Módulo re}

\begin{frame}[fragile]{Módulo re: Expresiones regulares en Python \faPython}
\begin{conceptbox}{¿Qué es el módulo re?}
Módulo de la librería estándar para trabajar con expresiones regulares
\end{conceptbox}

\textbf{Importación:}
\begin{lstlisting}
import re
\end{lstlisting}

\vspace{0.2cm}
\textbf{Funciones principales:}
\small
\begin{itemize}
    \item[\faSearch] \texttt{search()}: Busca patrón en cualquier parte del texto
    \item[\faAlignLeft] \texttt{match()}: Busca patrón al inicio del texto
    \item[\faListUl] \texttt{findall()}: Encuentra todas las coincidencias
    \item[\faExchangeAlt] \texttt{sub()}: Reemplaza coincidencias
    \item[\faCut] \texttt{split()}: Divide texto por patrón
\end{itemize}
\end{frame}

\section{Función search()}

\begin{frame}[fragile]{search(): Buscar patrón \faSearch}
\begin{conceptbox}{¿Qué hace?}
Busca la primera coincidencia del patrón en cualquier parte del texto
\end{conceptbox}

\begin{lstlisting}
import re

texto = "Mi teléfono es 555-1234"
patron = r"\d+"  # Uno o más dígitos

resultado = re.search(patron, texto)

if resultado:
    print(f"Encontrado: {resultado.group()}")
    print(f"Posición: {resultado.start()}-{resultado.end()}")
else:
    print("No se encontró")
\end{lstlisting}
\end{frame}

\begin{frame}[fragile]{search(): Salida \faArrowCircleRight}
\begin{tcolorbox}[colback=green!10,colframe=green,title={\small\faArrowCircleRight \ Salida},coltitle=fg0,left=3pt,right=3pt,top=3pt,bottom=3pt,boxsep=2pt]
\small\texttt{Encontrado: 555\\
Posición: 16-19}
\end{tcolorbox}

\vspace{0.3cm}
\begin{notebox}
\small\faInfoCircle\ La \texttt{r} antes del string indica "raw string" (importante para regex)
\end{notebox}
\end{frame}

\section{Función match()}

\begin{frame}[fragile]{match(): Buscar al inicio \faAlignLeft}
\begin{conceptbox}{¿Qué hace?}
Busca el patrón solo al inicio del texto
\end{conceptbox}

\begin{lstlisting}
import re

texto1 = "Python es genial"
texto2 = "Me gusta Python"

patron = r"Python"

match1 = re.match(patron, texto1)
match2 = re.match(patron, texto2)

print(f"Texto 1: {match1 is not None}")
print(f"Texto 2: {match2 is not None}")
\end{lstlisting}
\end{frame}

\begin{frame}[fragile]{match(): Salida \faArrowCircleRight}
\begin{tcolorbox}[colback=green!10,colframe=green,title={\small\faArrowCircleRight \ Salida},coltitle=fg0,left=3pt,right=3pt,top=3pt,bottom=3pt,boxsep=2pt]
\small\texttt{Texto 1: True\\
Texto 2: False}
\end{tcolorbox}

\vspace{0.3cm}
\begin{notebox}
\small\faExclamationTriangle\ \texttt{match()} solo busca al inicio, \texttt{search()} busca en cualquier parte
\end{notebox}
\end{frame}

\section{Función findall()}

\begin{frame}[fragile]{findall(): Encontrar todas las coincidencias \faListUl}
\begin{conceptbox}{¿Qué hace?}
Retorna una lista con todas las coincidencias del patrón
\end{conceptbox}

\begin{lstlisting}
import re

texto = "Precios: $10, $25, $100 y $5"
patron = r"\$\d+"  # $ seguido de dígitos

precios = re.findall(patron, texto)

print(f"Precios encontrados: {precios}")
print(f"Total: {len(precios)} precios")
\end{lstlisting}
\end{frame}

\begin{frame}[fragile]{findall(): Salida \faArrowCircleRight}
\begin{tcolorbox}[colback=green!10,colframe=green,title={\small\faArrowCircleRight \ Salida},coltitle=fg0,left=3pt,right=3pt,top=3pt,bottom=3pt,boxsep=2pt]
\small\texttt{Precios encontrados: ['\$10', '\$25', '\$100', '\$5']\\
Total: 4 precios}
\end{tcolorbox}

\vspace{0.3cm}
\begin{notebox}
\small\faInfoCircle\ \texttt{\textbackslash\$} escapa el carácter especial \$
\end{notebox}
\end{frame}

\section{Función sub()}

\begin{frame}[fragile]{sub(): Reemplazar texto \faExchangeAlt}
\begin{conceptbox}{¿Qué hace?}
Reemplaza todas las coincidencias del patrón por otro texto
\end{conceptbox}

\begin{lstlisting}
import re

texto = "Me gustan los gatos. Los gatos son geniales."
patron = r"gatos"
reemplazo = "perros"

nuevo_texto = re.sub(patron, reemplazo, texto)

print("Original:")
print(texto)
print("Modificado:")
print(nuevo_texto)
\end{lstlisting}
\end{frame}

\begin{frame}[fragile]{sub(): Salida \faArrowCircleRight}
\begin{tcolorbox}[colback=green!10,colframe=green,title={\small\faArrowCircleRight \ Salida},coltitle=fg0,left=3pt,right=3pt,top=3pt,bottom=3pt,boxsep=2pt]
\small\texttt{Original:\\
Me gustan los gatos. Los gatos son geniales.\\
Modificado:\\
Me gustan los perros. Los perros son geniales.}
\end{tcolorbox}

\vspace{0.3cm}
\begin{notebox}
\small\faCheckCircle\ \texttt{sub()} reemplaza todas las coincidencias
\end{notebox}
\end{frame}

\begin{frame}[fragile]{sub(): Limitar reemplazos \faEdit}
\begin{lstlisting}
import re

texto = "Ana tiene 3 gatos, Pedro tiene 2 gatos"
patron = r"\d+"
reemplazo = "X"

# Reemplazar solo la primera coincidencia
nuevo_texto = re.sub(patron, reemplazo, texto, count=1)

print(f"Original: {texto}")
print(f"Modificado: {nuevo_texto}")
\end{lstlisting}

\vspace{0.3cm}
\begin{tcolorbox}[colback=green!10,colframe=green,title={\small\faArrowCircleRight \ Salida},coltitle=fg0,left=3pt,right=3pt,top=3pt,bottom=3pt,boxsep=2pt]
\small\texttt{Original: Ana tiene 3 gatos, Pedro tiene 2 gatos\\
Modificado: Ana tiene X gatos, Pedro tiene 2 gatos}
\end{tcolorbox}
\end{frame}

\begin{frame}{sub(): Nota sobre count \faLightbulb}
\begin{notebox}
\small\faLightbulb\ El parámetro \texttt{count} limita el número de reemplazos a realizar
\end{notebox}
\end{frame}

\section{Patrones básicos}

\begin{frame}[fragile]{Patrones literales \faFont}
\begin{conceptbox}{Caracteres literales}
Los caracteres normales buscan coincidencias exactas
\end{conceptbox}

\begin{lstlisting}
import re

texto = "Python es Python, y Python es genial"

# Buscar "Python" literalmente
coincidencias = re.findall(r"Python", texto)
print(f"Coincidencias: {len(coincidencias)}")

# Buscar "python" (mayúsculas/minúsculas importan)
coincidencias2 = re.findall(r"python", texto)
print(f"Con minúscula: {len(coincidencias2)}")
\end{lstlisting}
\end{frame}

\begin{frame}[fragile]{Patrones literales: Salida \faArrowCircleRight}
\begin{tcolorbox}[colback=green!10,colframe=green,title={\small\faArrowCircleRight \ Salida},coltitle=fg0,left=3pt,right=3pt,top=3pt,bottom=3pt,boxsep=2pt]
\small\texttt{Coincidencias: 3\\
Con minúscula: 0}
\end{tcolorbox}

\vspace{0.3cm}
\begin{notebox}
\small\faExclamationTriangle\ Por defecto, regex distingue mayúsculas/minúsculas
\end{notebox}
\end{frame}

\section{Metacaracteres}

\begin{frame}{Metacaracteres especiales \faMagic}
\begin{conceptbox}{Caracteres con significado especial}
\small
\begin{tabular}{ll}
\texttt{.} & Cualquier carácter (excepto salto de línea) \\
\texttt{\^} & Inicio del texto \\
\texttt{\$} & Final del texto \\
\texttt{|} & O (alternativa) \\
\texttt{[]} & Conjunto de caracteres \\
\texttt{()} & Grupo de captura \\
\texttt{\textbackslash} & Escape de caracteres especiales \\
\end{tabular}
\end{conceptbox}

\vspace{0.2cm}
\begin{center}
\textcolor{brown}{\faInfoCircle\ Estos caracteres tienen funciones especiales}
\end{center}
\end{frame}

\begin{frame}[fragile]{Metacaracteres en acción \faPlayCircle}
\begin{lstlisting}
import re

texto = "gato, pato, plato, rato"

# . = cualquier carácter
patron1 = r".ato"
print(f"Patrón '.ato': {re.findall(patron1, texto)}")

# ^ = inicio
patron2 = r"^gato"
print(f"Empieza con 'gato': {re.search(patron2, texto) is not None}")

# $ = final
patron3 = r"rato$"
print(f"Termina con 'rato': {re.search(patron3, texto) is not None}")
\end{lstlisting}
\end{frame}

\begin{frame}[fragile]{Metacaracteres: Salida \faArrowCircleRight}
\begin{tcolorbox}[colback=green!10,colframe=green,title={\small\faArrowCircleRight \ Salida},coltitle=fg0,left=3pt,right=3pt,top=3pt,bottom=3pt,boxsep=2pt]
\small\texttt{Patrón '.ato': ['gato', 'pato', 'lato', 'rato']\\
Empieza con 'gato': True\\
Termina con 'rato': True}
\end{tcolorbox}
\end{frame}

\begin{frame}[fragile]{Conjuntos de caracteres [] \faSquare}
\begin{lstlisting}
import re

texto = "a1 b2 c3 d4 e5"

# [abc] = a, b o c
patron1 = r"[abc]\d"
print(f"[abc] + dígito: {re.findall(patron1, texto)}")

# [a-z] = cualquier letra minúscula
patron2 = r"[a-z]\d"
print(f"Letra + dígito: {re.findall(patron2, texto)}")

# [^0-9] = cualquier cosa excepto dígitos
patron3 = r"[^0-9]+"
print(f"No dígitos: {re.findall(patron3, texto)}")
\end{lstlisting}
\end{frame}

\begin{frame}[fragile]{Conjuntos: Salida \faArrowCircleRight}
\begin{tcolorbox}[colback=green!10,colframe=green,title={\small\faArrowCircleRight \ Salida},coltitle=fg0,left=3pt,right=3pt,top=3pt,bottom=3pt,boxsep=2pt]
\small\texttt{[abc] + dígito: ['a1', 'b2', 'c3']\\
Letra + dígito: ['a1', 'b2', 'c3', 'd4', 'e5']\\
No dígitos: ['a', ' b', ' c', ' d', ' e']}
\end{tcolorbox}
\end{frame}

\section{Clases de caracteres}

\begin{frame}{Clases de caracteres predefinidas \faKeyboard}
\begin{conceptbox}{Atajos útiles}
\small
\begin{tabular}{ll}
\texttt{\textbackslash d} & Dígito: [0-9] \\
\texttt{\textbackslash D} & No dígito: [\^{}0-9] \\
\texttt{\textbackslash w} & Palabra: [a-zA-Z0-9\_] \\
\texttt{\textbackslash W} & No palabra \\
\texttt{\textbackslash s} & Espacio en blanco (espacio, tab, salto) \\
\texttt{\textbackslash S} & No espacio en blanco \\
\end{tabular}
\end{conceptbox}

\vspace{0.2cm}
\begin{center}
\textcolor{brown}{\faLightbulb\ Minúscula incluye, MAYÚSCULA excluye}
\end{center}
\end{frame}

\begin{frame}[fragile]{Clases de caracteres en acción \faCode}
\begin{lstlisting}
import re

texto = "Hola123 mundo_456 ¡Python!"

# \d = dígitos
digitos = re.findall(r"\d+", texto)
print(f"Dígitos: {digitos}")

# \w = caracteres de palabra
palabras = re.findall(r"\w+", texto)
print(f"Palabras: {palabras}")

# \s = espacios
espacios = re.findall(r"\s", texto)
print(f"Espacios: {len(espacios)}")
\end{lstlisting}
\end{frame}

\begin{frame}[fragile]{Clases: Salida \faArrowCircleRight}
\begin{tcolorbox}[colback=green!10,colframe=green,title={\small\faArrowCircleRight \ Salida},coltitle=fg0,left=3pt,right=3pt,top=3pt,bottom=3pt,boxsep=2pt]
\small\texttt{Dígitos: ['123', '456']\\
Palabras: ['Hola123', 'mundo\_456', 'Python']\\
Espacios: 2}
\end{tcolorbox}

\vspace{0.3cm}
\begin{notebox}
\small\faInfoCircle\ \texttt{\textbackslash w} no incluye caracteres especiales como ¡
\end{notebox}
\end{frame}

\section{Cuantificadores}

\begin{frame}{Cuantificadores: Repeticiones \faRedoAlt}
\begin{conceptbox}{Especifican cuántas veces debe aparecer}
\small
\begin{tabular}{ll}
\texttt{*} & 0 o más veces \\
\texttt{+} & 1 o más veces \\
\texttt{?} & 0 o 1 vez (opcional) \\
\texttt{\{n\}} & Exactamente n veces \\
\texttt{\{n,\}} & n o más veces \\
\texttt{\{n,m\}} & Entre n y m veces \\
\end{tabular}
\end{conceptbox}

\vspace{0.2cm}
\begin{center}
\textcolor{brown}{\faExclamationTriangle\ Los cuantificadores son "codiciosos" por defecto}
\end{center}
\end{frame}

\begin{frame}[fragile]{Cuantificadores básicos: *, +, ? \faAsterisk}
\begin{lstlisting}
import re

texto = "Color, Colour, Colouur, Colr"

# o* = cero o más 'o'
patron1 = r"Colo*r"
print(f"o*: {re.findall(patron1, texto)}")

# o+ = una o más 'o'
patron2 = r"Colo+r"
print(f"o+: {re.findall(patron2, texto)}")

# o? = cero o una 'o'
patron3 = r"Colo?r"
print(f"o?: {re.findall(patron3, texto)}")
\end{lstlisting}
\end{frame}

\begin{frame}[fragile]{Cuantificadores: Salida \faArrowCircleRight}
\begin{tcolorbox}[colback=green!10,colframe=green,title={\small\faArrowCircleRight \ Salida},coltitle=fg0,left=3pt,right=3pt,top=3pt,bottom=3pt,boxsep=2pt]
\small\texttt{o*: ['Color', 'Colour', 'Colouur', 'Colr']\\
o+: ['Color', 'Colour', 'Colouur']\\
o?: ['Color', 'Colr']}
\end{tcolorbox}
\end{frame}

\begin{frame}[fragile]{Cuantificadores exactos: \{n,m\} \faHashtag}
\begin{lstlisting}
import re

texto = "Números: 1, 12, 123, 1234, 12345"

# Exactamente 3 dígitos
patron1 = r"\d{3}"
print(f"3 dígitos: {re.findall(patron1, texto)}")

# Entre 2 y 4 dígitos
patron2 = r"\d{2,4}"
print(f"2-4 dígitos: {re.findall(patron2, texto)}")

# 2 o más dígitos
patron3 = r"\d{2,}"
print(f"2+ dígitos: {re.findall(patron3, texto)}")
\end{lstlisting}
\end{frame}

\begin{frame}[fragile]{Cuantificadores exactos: Salida \faArrowCircleRight}
\begin{tcolorbox}[colback=green!10,colframe=green,title={\small\faArrowCircleRight \ Salida},coltitle=fg0,left=3pt,right=3pt,top=3pt,bottom=3pt,boxsep=2pt]
\small\texttt{3 dígitos: ['123', '123', '123']\\
2-4 dígitos: ['12', '123', '1234', '1234']\\
2+ dígitos: ['12', '123', '1234', '12345']}
\end{tcolorbox}

\vspace{0.3cm}
\begin{notebox}
\small\faInfoCircle\ Los cuantificadores coinciden lo máximo posible
\end{notebox}
\end{frame}

\section{Grupos de captura}

\begin{frame}[fragile]{Grupos de captura () \faObjectGroup}
\begin{conceptbox}{¿Para qué sirven?}
Capturan partes específicas del patrón para extraerlas o referenciarlas
\end{conceptbox}

\begin{lstlisting}
import re

texto = "Fecha: 2024-10-21"

# Capturar año, mes y día por separado
patron = r"(\d{4})-(\d{2})-(\d{2})"

match = re.search(patron, texto)

if match:
    print(f"Fecha completa: {match.group(0)}")
    print(f"Año: {match.group(1)}")
    print(f"Mes: {match.group(2)}")
    print(f"Día: {match.group(3)}")
\end{lstlisting}
\end{frame}

\begin{frame}[fragile]{Grupos de captura: Salida \faArrowCircleRight}
\begin{tcolorbox}[colback=green!10,colframe=green,title={\small\faArrowCircleRight \ Salida},coltitle=fg0,left=3pt,right=3pt,top=3pt,bottom=3pt,boxsep=2pt]
\small\texttt{Fecha completa: 2024-10-21\\
Año: 2024\\
Mes: 10\\
Día: 21}
\end{tcolorbox}
\end{frame}

\begin{frame}[fragile]{Grupos nombrados (?P<nombre>) \faTag}
\begin{lstlisting}
import re

texto = "Email: usuario@dominio.com"

# Grupos nombrados para mayor claridad
patron = r"(?P<user>\w+)@(?P<domain>\w+\.\w+)"

match = re.search(patron, texto)

if match:
    print(f"Usuario: {match.group('user')}")
    print(f"Dominio: {match.group('domain')}")
    print(f"Dict: {match.groupdict()}")
\end{lstlisting}
\end{frame}

\begin{frame}[fragile]{Grupos nombrados: Salida \faArrowCircleRight}
\begin{tcolorbox}[colback=green!10,colframe=green,title={\small\faArrowCircleRight \ Salida},coltitle=fg0,left=3pt,right=3pt,top=3pt,bottom=3pt,boxsep=2pt]
\small\texttt{Usuario: usuario\\
Dominio: dominio.com\\
Dict: \{'user': 'usuario', 'domain': 'dominio.com'\}}
\end{tcolorbox}

\vspace{0.3cm}
\begin{notebox}
\small\faLightbulb\ Los grupos nombrados son más legibles y mantenibles
\end{notebox}
\end{frame}

\section{Validación de formatos}

\begin{frame}[fragile]{Validar emails \faEnvelope}
\begin{lstlisting}
import re

def validar_email(email):
    patron = r"^[a-zA-Z0-9._%+-]+@[a-zA-Z0-9.-]+\.[a-zA-Z]{2,}$"
    return re.match(patron, email) is not None

emails = ["usuario@ejemplo.com", "nombre.apellido@empresa.co",
          "correo_invalido@", "@sin-usuario.com"]

for email in emails:
    valido = validar_email(email)
    simbolo = "OK" if valido else "X"
    print(f"{email:28} -> {simbolo}")
\end{lstlisting}
\end{frame}

\begin{frame}[fragile]{Validar emails: Salida \faArrowCircleRight}
\begin{tcolorbox}[colback=green!10,colframe=green,title={\small\faArrowCircleRight \ Salida},coltitle=fg0,left=3pt,right=3pt,top=3pt,bottom=3pt,boxsep=2pt]
\small\texttt{usuario@ejemplo.com        -> OK\\
nombre.apellido@empresa.co -> OK\\
correo\_invalido@           -> X\\
@sin-usuario.com           -> X}
\end{tcolorbox}
\end{frame}

\begin{frame}[fragile]{Validar teléfonos \faPhone}
\begin{lstlisting}
import re

def validar_telefono(telefono):
    patrones = [
        r"^\d{3}-\d{3}-\d{4}$",
        r"^\(\d{3}\)\s\d{3}-\d{4}$",
        r"^\d{10}$"
    ]
    for patron in patrones:
        if re.match(patron, telefono):
            return True
    return False

tels = ["555-123-4567", "(555) 123-4567", "5551234567", "123"]
for tel in tels:
    print(f"{tel:18} -> {'OK' if validar_telefono(tel) else 'X'}")
\end{lstlisting}
\end{frame}

\begin{frame}[fragile]{Validar teléfonos: Salida \faArrowCircleRight}
\begin{tcolorbox}[colback=green!10,colframe=green,title={\small\faArrowCircleRight \ Salida},coltitle=fg0,left=3pt,right=3pt,top=3pt,bottom=3pt,boxsep=2pt]
\small\texttt{555-123-4567       -> OK\\
(555) 123-4567     -> OK\\
5551234567         -> OK\\
123                -> X}
\end{tcolorbox}
\end{frame}

\begin{frame}[fragile]{Validar URLs \faLink}
\begin{lstlisting}
import re

def validar_url(url):
    patron = r"^https?://[a-zA-Z0-9.-]+\.[a-zA-Z]{2,}(/.*)?$"
    return re.match(patron, url) is not None

urls = ["https://www.ejemplo.com",
        "http://ejemplo.com/pagina",
        "https://sub.dominio.com/ruta/archivo.html",
        "www.sinhttp.com",
        "ftp://protocolo-invalido.com"]

for url in urls:
    resultado = "OK" if validar_url(url) else "X"
    print(f"{resultado:2} {url[:35]}")
\end{lstlisting}
\end{frame}

\begin{frame}[fragile]{Validar URLs: Salida \faArrowCircleRight}
\begin{tcolorbox}[colback=green!10,colframe=green,title={\small\faArrowCircleRight \ Salida},coltitle=fg0,left=3pt,right=3pt,top=3pt,bottom=3pt,boxsep=2pt]
\small\texttt{OK https://www.ejemplo.com\\
OK http://ejemplo.com/pagina\\
OK https://sub.dominio.com/ruta/archi\\
X  www.sinhttp.com\\
X  ftp://protocolo-invalido.com}
\end{tcolorbox}
\end{frame}

\section{Banderas (flags)}

\begin{frame}[fragile]{Banderas: Modificar comportamiento \faFlag}
\begin{conceptbox}{Flags comunes}
\small
\begin{itemize}
    \item[\faFont] \texttt{re.IGNORECASE} o \texttt{re.I}: Ignora mayúsculas
    \item[\faAlignJustify] \texttt{re.MULTILINE} o \texttt{re.M}: \^{} y \$ por línea
    \item[\faEllipsisH] \texttt{re.DOTALL} o \texttt{re.S}: . incluye saltos
\end{itemize}
\end{conceptbox}

\begin{lstlisting}
import re

texto = "Python es PYTHON y python"

# Sin bandera
sin_flag = re.findall(r"python", texto)
print(f"Sin flag: {sin_flag}")

# Con IGNORECASE
con_flag = re.findall(r"python", texto, re.IGNORECASE)
print(f"Con IGNORECASE: {con_flag}")
\end{lstlisting}
\end{frame}

\begin{frame}[fragile]{Banderas: Salida \faArrowCircleRight}
\begin{tcolorbox}[colback=green!10,colframe=green,title={\small\faArrowCircleRight \ Salida},coltitle=fg0,left=3pt,right=3pt,top=3pt,bottom=3pt,boxsep=2pt]
\small\texttt{Sin flag: ['python']\\
Con IGNORECASE: ['Python', 'PYTHON', 'python']}
\end{tcolorbox}
\end{frame}

\section{Compilar patrones}

\begin{frame}[fragile]{Compilar patrones para reutilizar \faSync}
\begin{conceptbox}{¿Por qué compilar?}
Mejora el rendimiento cuando usas el mismo patrón varias veces
\end{conceptbox}

\begin{lstlisting}
import re

patron_email = re.compile(
    r"[a-zA-Z0-9._%+-]+@[a-zA-Z0-9.-]+\.[a-zA-Z]{2,}")

textos = ["Contacto: juan@ejemplo.com",
          "Soporte: ayuda@empresa.com",
          "Sin email aquí"]

for texto in textos:
    match = patron_email.search(texto)
    if match:
        print(f"Email: {match.group()}")
    else:
        print("No se encontró email")
\end{lstlisting}
\end{frame}

\begin{frame}[fragile]{Compilar patrones: Salida \faArrowCircleRight}
\begin{tcolorbox}[colback=green!10,colframe=green,title={\small\faArrowCircleRight \ Salida},coltitle=fg0,left=3pt,right=3pt,top=3pt,bottom=3pt,boxsep=2pt]
\small\texttt{Email: juan@ejemplo.com\\
Email: ayuda@empresa.com\\
No se encontró email}
\end{tcolorbox}
\end{frame}

\section{Ejemplo integrador}

\begin{frame}[fragile]{Proyecto: Extractor de información (1/5) \faFileAlt}
\textbf{Objetivo:} Extraer información de un texto con múltiples formatos

\vspace{0.2cm}
\textbf{Archivo:} \texttt{extractor.py}
\begin{lstlisting}
import re

class ExtractorInfo:
    """Extrae emails, teléfonos, URLs y fechas de texto"""
    
    def __init__(self):
        # Compilar patrones
        self.patron_email = re.compile(
            r"[a-zA-Z0-9._%+-]+@[a-zA-Z0-9.-]+\.[a-zA-Z]{2,}"
        )
        self.patron_telefono = re.compile(
            r"\(?\d{3}\)?[-.\s]?\d{3}[-.\s]?\d{4}"
        )
        self.patron_url = re.compile(
            r"https?://[a-zA-Z0-9.-]+\.[a-zA-Z]{2,}(/[^\s]*)?"
        )
\end{lstlisting}
\end{frame}

\begin{frame}[fragile]{Proyecto: Extractor de información (2/5) \faCode}
\begin{lstlisting}
        self.patron_fecha = re.compile(
            r"\d{1,2}[-/]\d{1,2}[-/]\d{2,4}"
        )
    
    def extraer_emails(self, texto):
        """Extrae todos los emails del texto"""
        return self.patron_email.findall(texto)
    
    def extraer_telefonos(self, texto):
        """Extrae todos los teléfonos del texto"""
        return self.patron_telefono.findall(texto)
    
    def extraer_urls(self, texto):
        """Extrae todas las URLs del texto"""
        return self.patron_url.findall(texto)
    
    def extraer_fechas(self, texto):
        """Extrae todas las fechas del texto"""
        return self.patron_fecha.findall(texto)
\end{lstlisting}
\end{frame}

\begin{frame}[fragile]{Proyecto: Extractor de información (3/5) \faSearch}
\begin{lstlisting}
    def extraer_todo(self, texto):
        """Extrae toda la información del texto"""
        return {
            "emails": self.extraer_emails(texto),
            "telefonos": self.extraer_telefonos(texto),
            "urls": self.extraer_urls(texto),
            "fechas": self.extraer_fechas(texto)
        }
    
    def mostrar_resultados(self, resultados):
        """Muestra los resultados de forma legible"""
        print("\n=== INFORMACIÓN EXTRAÍDA ===\n")
        
        for categoria, items in resultados.items():
            print(f"{categoria.upper()}:")
            if items:
                for item in items:
                    print(f"  - {item}")
            else:
                print("  (ninguno)")
            print()
\end{lstlisting}
\end{frame}

\begin{frame}[fragile]{Proyecto: Extractor de información (4/5) \faPlayCircle}
\begin{lstlisting}
texto_ejemplo = """
Hola, mi email es contacto@ejemplo.com y también
puedes escribirme a usuario123@empresa.co

Llámame al (555) 123-4567 o al 555-987-6543

Visita: https://www.ejemplo.com
Más info: http://docs.ejemplo.com/guia

Reunión: 21-10-2024
Siguiente cita: 15/11/2024
"""

extractor = ExtractorInfo()
resultados = extractor.extraer_todo(texto_ejemplo)
extractor.mostrar_resultados(resultados)
\end{lstlisting}
\end{frame}

\begin{frame}[fragile]{Proyecto: Extractor de información (5/5) \faTerminal}
\begin{tcolorbox}[colback=green!10,colframe=green,title={\small\faArrowCircleRight \ Salida},coltitle=fg0,left=3pt,right=3pt,top=2pt,bottom=2pt,boxsep=1pt]
\tiny\texttt{=== INFORMACIÓN EXTRAÍDA ===\\
\ \\
EMAILS:\\
\ \ - contacto@ejemplo.com\\
\ \ - usuario123@empresa.co\\
\ \\
TELEFONOS:\\
\ \ - (555) 123-4567\\
\ \ - 555-987-6543\\
\ \\
URLS:\\
\ \ - https://www.ejemplo.com\\
\ \ - http://docs.ejemplo.com/guia\\
\ \\
FECHAS:\\
\ \ - 21-10-2024\\
\ \ - 15/11/2024}
\end{tcolorbox}

\vspace{-0.1cm}
\begin{notebox}
\small\faRocket\ ¡Extractor completo y reutilizable!
\end{notebox}
\end{frame}

\begin{frame}[fragile]{Proyecto alternativo: Validador de formularios \faCheckSquare}
\begin{lstlisting}
import re

class ValidadorFormulario:
    
    @staticmethod
    def validar_nombre(nombre):
        patron = r"^[a-zA-ZáéíóúÁÉÍÓÚñÑ\s]{2,50}$"
        return re.match(patron, nombre) is not None
    
    @staticmethod
    def validar_email(email):
        p = r"^[a-zA-Z0-9._%+-]+@[a-zA-Z0-9.-]+\.[a-zA-Z]{2,}$"
        return re.match(p, email) is not None
\end{lstlisting}

\small(Continuación en siguiente slide...)
\end{frame}

\begin{frame}[fragile]{Proyecto alternativo: Validador (cont.) \faEdit}
\begin{lstlisting}
    @staticmethod
    def validar_password(password):
        if len(password) < 8:
            return False
        tiene_mayus = re.search(r"[A-Z]", password)
        tiene_minus = re.search(r"[a-z]", password)
        tiene_num = re.search(r"\d", password)
        return all([tiene_mayus, tiene_minus, tiene_num])
    
    def validar_formulario(self, datos):
        res = {}
        res["nombre"] = self.validar_nombre(datos["nombre"])
        res["email"] = self.validar_email(datos["email"])
        res["password"] = self.validar_password(datos["password"])
        return res
\end{lstlisting}
\end{frame}

\begin{frame}[fragile]{Proyecto alternativo: Uso del validador \faTerminal}
\begin{lstlisting}
# Probar validador
validador = ValidadorFormulario()

formulario = {
    "nombre": "Juan Pérez",
    "email": "juan@ejemplo.com",
    "password": "MiClave123"
}

resultados = validador.validar_formulario(formulario)

print("=== VALIDACIÓN DE FORMULARIO ===")
for campo, valido in resultados.items():
    estado = "✓ Válido" if valido else "✗ Inválido"
    print(f"{campo:10} -> {estado}")
\end{lstlisting}
\end{frame}

\begin{frame}[fragile]{Proyecto alternativo: Salida \faArrowCircleRight}
\begin{tcolorbox}[colback=green!10,colframe=green,title={\small\faArrowCircleRight \ Salida},coltitle=fg0,left=3pt,right=3pt,top=3pt,bottom=3pt,boxsep=2pt]
\small\texttt{=== VALIDACIÓN DE FORMULARIO ===\\
nombre     -> ✓ Válido\\
email      -> ✓ Válido\\
password   -> ✓ Válido}
\end{tcolorbox}
\end{frame}

\begin{frame}{Buenas prácticas con regex \faCheckDouble}
\begin{conceptbox}{Recomendaciones}
\small
\begin{itemize}
    \item[\faCheck] Usa raw strings (\texttt{r""}) para patrones
    \item[\faCheck] Compila patrones si los reutilizas
    \item[\faCheck] Nombra grupos de captura para claridad
    \item[\faCheck] Comenta patrones complejos
    \item[\faCheck] Prueba patrones con varios casos
    \item[\faCheck] No uses regex para parsear HTML/XML
    \item[\faCheck] Mantén patrones simples cuando sea posible
    \item[\faCheck] Usa herramientas online para probar (regex101.com)
\end{itemize}
\end{conceptbox}
\end{frame}

\begin{frame}[fragile]{Herramientas útiles \faTools}
\begin{conceptbox}{Recursos para trabajar con regex}
\small
\begin{itemize}
    \item[\faGlobe] \textbf{regex101.com}: Probador online interactivo
    \item[\faGlobe] \textbf{regexr.com}: Aprendizaje visual
    \item[\faBook] \textbf{Python docs}: Documentación oficial
    \item[\faTerminal] \textbf{re.DEBUG}: Análisis del patrón
\end{itemize}
\end{conceptbox}

\begin{lstlisting}
import re

# Ver análisis interno del patrón
patron = re.compile(r"\d{3}-\d{4}", re.DEBUG)
\end{lstlisting}
\end{frame}

\begin{frame}{Herramientas: Nota \faLightbulb}
\begin{notebox}
\small\faLightbulb\ Regex puede ser complejo, ¡practica y usa herramientas!
\end{notebox}
\end{frame}

\begin{frame}{Resumen \faListUl}
\begin{conceptbox}{¿Qué aprendimos hoy?}
\small
\begin{itemize}
    \item[\faCheck] Introducción a expresiones regulares
    \item[\faCheck] Módulo \texttt{re}: search, match, findall, sub
    \item[\faCheck] Patrones básicos y metacaracteres
    \item[\faCheck] Clases: \textbackslash d, \textbackslash w, \textbackslash s
    \item[\faCheck] Cuantificadores: *, +, ?, \{n,m\}
    \item[\faCheck] Grupos de captura y nombrados
    \item[\faCheck] Validación de formatos
    \item[\faCheck] Proyecto: Extractor de información
\end{itemize}
\end{conceptbox}

\vspace{0.3cm}
\begin{center}
\large\textcolor{brown}{\faSearch\ ¡Domina el texto con regex!}
\end{center}
\end{frame}

\begin{frame}
\begin{center}
\Huge\textcolor{brown}{\faQuestionCircle}\\
\vspace{0.5cm}
\Huge ¡Preguntas!\\
\vspace{1cm}
\Large\textcolor{brownbright}{Gracias por su atención}
\end{center}
\end{frame}

\end{document}
